Next we turn to investigating the content of firms' communications. For the \data{audit/true_positive_count_int} true positive transcripts, the Supplementary Information categorizes them according to the type of collusive content. Here, we focus on the five most common categories which are:\footnote{Some excerpts have content spanning more than one category, which will be apparent in some of the quoted passages. When the entire transcript is read, there are likely to be even more cases in which content encompassed multiple categories. For this reason, the number assigned to each category is not meant to be definitive but only to give a rough sense of its frequency.}
\begin{itemize}
    \item A firm announces the industry has reduced supply or capacity and references competitors or the industry in a manner supportive of such conduct. (40 transcripts)
    \item A firm announces the industry needs to reduce supply or capacity or to show discipline. (30 transcripts)
    \item A firm announces the industry has shown or will show supply or capacity or price discipline. (12 transcripts)
    \item A firm announces the industry needs to raise prices. (12 transcripts)
    \item A firm announces it is being or will be less aggressive in pricing and references competitors’ conduct in a manner encouraging similar conduct. (8 transcripts)
\end{itemize}

\subsubsection{Category 1: The industry has reduced supply or capacity and references competitors or the industry in a manner supportive of such conduct}
The most common collusive content is that the industry has reduced supply or capacity and refers to competitors or the industry to facilitate such conduct. This type of content is present in 37\% of the 109 true positive transcripts. In emphasizing supply or capacity cuts, a firm may highlight how those reductions are in the industry’s best interests. For example:\footnote{For each quoted excerpt, we include in square brackets the company name, event type, and date, as reported in Compustat, corresponding to the original transcript that was provided in the LLM query. While previously documented cases like airlines are in the set of true positives, we chose not to use excerpts from them (even though some are highly illustrative) in order to focus on new episodes identified by the LLM. As the passages are from the LLM-generated excerpts, the sentences may be from different parts of the transcript and not appear together as they are presented here.}

\begin{quote}
``We've seen a massive … consolidation in the industry ... and hopefully, that leads to a better, a healthier competitive environment as well. No one is going to build DRAM capacity for the foreseeable future. It's just not going to happen.'' \\
\mbox{[\data{quoted_excerpt_labels/audit_types/q508027}]}
\end{quote}

\begin{quote}
``Other companies in our industry have announced considerable domestic downtime and capacity closures [and] share my philosophy that if you can't sell it or you can't use it, you shouldn't make it.'' \\
\mbox{[\data{quoted_excerpt_labels/audit_types/q597966}]}
\end{quote}

\begin{quote}
``We will not continue to invest at this pace, and our activity will slow. It's kind of a no-brainer from an economic standpoint. Our partners are making the right decisions in terms of the decreasing capital in this environment.'' \\
\mbox{[\data{quoted_excerpt_labels/audit_types/q419393}]}
\end{quote}

In some cases, the firm went beyond noting a decline in supply or capacity to express more reductions are needed.

\begin{quote}
``We currently anticipate these conditions will require the industry to rationalize production to stabilize the markets so fundamentals can improve. While some supply response has already occurred, more reductions are necessary to balance the market.'' \\
\mbox{[\data{quoted_excerpt_labels/audit_types/q185218}]}
\end{quote}

\begin{quote}
``On the market side, [there is] still [a] challenging environment in Europe. Industry-wide capacity closures are supportive [and] if we can see maybe 1 or 2 more machine 
closures in Europe, I think then we are more or less moving it.'' \\
\mbox{[\data{quoted_excerpt_labels/audit_types/q260640}]}
\end{quote}

\begin{quote}
``There continues to be too much capacity in the global preclinical services industry, and the excess capacity keeps pressure on pricing. We recognized that we would have to reduce capacity in order to improve utilization. A competitor has just announced capacity reductions, and we expect that the pharma industry will continue its efforts to sell and/or close preclinical space. We believe that more capacity remains to be rationalized [and] these actions will help to improve the pricing dynamics.'' \\
\mbox{[\data{quoted_excerpt_labels/audit_types/q389910}]}
\end{quote}

In nine of the 40 transcripts, reduced industry supply or capacity is mentioned along with references to competitors doing so and the firm supporting such conduct.

\begin{quote}
``We are limiting our investments with reductions in capital [and] production declines will soon be evident. We are witnessing peer companies follow our lead [and] think more companies should adopt this prudent practice.'' \\
\mbox{[\data{quoted_excerpt_labels/audit_types/q419394}]}
\end{quote}

\begin{quote}
``Industry linerboard and Kraft paper capacity will be substantially reduced [and] we should expect to see further capacity closures announced in the coming months. [This] should lead to substantial price restoration [as] high operating rates give the industry courage to raise prices. I believe that's good for our shareholders and good for the industry.'' \\
\mbox{[\data{quoted_excerpt_labels/audit_types/q597969}]}
\end{quote}

\begin{quote}
``The industry is oversupplied in NAND, and we need to slow down bit supply. So everybody is slowing. I think the industry is reacting well. We're not trying to gain share. It's about margin profile and profitability.'' \\
\mbox{[\data{quoted_excerpt_labels/audit_types/q617313}]}
\end{quote}

And when competitors have not reduced supply or capacity, the firm calls them out for not doing so.

\begin{quote}
``We've also done what we think is the prudent thing to do [which] is to effectively provide ... incremental storage to the market by curtailing our production and not pulling it out of the ground. It has been a mystery to me why more producers have not made that same economic choice. We think that it has a chance to restore gas prices, but we'll see how that pans out.'' \\
\mbox{[\data{quoted_excerpt_labels/audit_types/q381391}]}
\end{quote}

In some instances, a firm notes that it is acting as a market leader in reducing supply or capacity. With that leadership reference comes the implicit invitation for competitors to be followers.

\begin{quote}
``We are currently idling temporarily ... approximately 40\% [of] operating capacity [and] we will be implementing a price increase of up to 10\%. There is a supply and demand imbalance in this business, and being a market leader, we have taken a step to reduce our capacity.'' \\
\mbox{[\data{quoted_excerpt_labels/audit_types/q261028}]}
\end{quote}

\begin{quote}
``We took our production level down [and] we believe our actions had a positive effect on bringing industry ethanol stocks into a more stable supply and demand balance. Somebody had to exert leadership in this industry to do it. And if we did not slow down, we would be in a much worse margin situation today.'' \\
\mbox{[\data{quoted_excerpt_labels/audit_types/q348864}]}
\end{quote}

\begin{quote}
``We have made production curtailments. Following our lead, we're now seeing a number of recent capacity curtailments by other producers as the industry takes action necessary to support a recovery in prices. We view this trend favorably and see this dynamic as critical to stabilizing and eventually improving the pricing environment.'' \\
\mbox{[\data{quoted_excerpt_labels/audit_types/q569063}]}
\end{quote}

\subsubsection{Category 2: The industry needs to reduce supply or capacity or to show discipline}
The second most common category is when a firm announces that the industry needs to reduce supply or capacity or to show discipline. These instances made up 28\% of the true positives. As documented in airlines and steel (see \textcite{aryal2022coordinated} and \textcite{harrington2022collusion}), ``discipline'' is a common code word for ``less competition''. Other loaded words include firms being ``responsible'' and ``rationalizing'' capacity.

\begin{quote}
``We do consider the possibility of reducing our production in 2017 in order to have a better market environment and also a better environment for pricing. If we in the industry … do not have a discipline regarding our supply, ... we could be generating a situation in which the industry's ROIC could be lower. We are not sure that this will bring other companies to follow in terms of the supply strategies. But what is clear to us is that [we have] the responsibility of showing the facts to be followed in order to seek profitability for the whole industry.'' \\
\mbox{[\data{quoted_excerpt_labels/audit_types/q406151}]}
\end{quote}

\begin{quote}
``It is important that the U.S. oil and gas industry as well as other global producers continue exerting capital discipline to not overproduce. Everyone needs to continue to cutback and participate as this slump continues. We reduced our pace of growth to deliver a disciplined approach to value creation.'' \\
\mbox{[\data{quoted_excerpt_labels/audit_types/q325885}]}
\end{quote}

\begin{quote}
``We plan to optimize near-term pricing through the voluntary curtailment of approximately 100 million cubic feet per day over the next several months. Our hope is that our peers, some by choice rather than necessity, will do the same. In today's market, I believe what our industry needs is a more measured pace of growth.'' \\
\mbox{[\data{quoted_excerpt_labels/audit_types/q228511}]}
\end{quote}

One firm went so far as to specify how much capacity needed to be cut.

\begin{quote}
``It's an obvious conclusion that scrapping has to take place [and] according to my calculations, about 1/3 of the global fleet should be scrapped. Consolidation is inevitable and will take place in our industry.'' \\
\mbox{[\data{quoted_excerpt_labels/audit_types/q494779}]}
\end{quote}

These expressions of an industry need to cut supply or capacity may be explained in the context of raising prices.

\begin{quote}
``The average prices over the last 15 years in Europe has been EUR 550 … and we believe that ... the trend prices should be more in this range than the historically very low levels that we've seen this year of EUR 420, EUR 430.'' \\
\mbox{[\data{quoted_excerpt_labels/audit_types/q260628}]}
\end{quote}

The firm may go on to note being encouraged by competitors’ conduct.

\begin{quote}
``The market has overcapacity now, which puts the leverage in the customers' hands. You can resolve overcapacity by not making investments. What needs to happen is capacity needs to shrink. I've been encouraged by what I see our competitors doing.'' \\
\mbox{[\data{quoted_excerpt_labels/audit_types/q199854}]}
\end{quote}

Especially egregious is when a firm announces how its own supply or capacity cuts reflect ``doing its part by exerting discipline'' [352688] and then demands competitors to pull their weight.

\begin{quote}
``We've implemented planned supply discipline. We think we've done our share and maybe more.'' \\
\mbox{[\data{quoted_excerpt_labels/audit_types/q519067}]}
\end{quote}

\begin{quote}
``And candidly, we can't do it alone. It's got to be an industry level. There's a gross oversupply in the space. We're taking aggressive action on the supply side.'' \\
\mbox{[\data{quoted_excerpt_labels/audit_types/q508184}]}
\end{quote}

\begin{quote}
``The market is clearly oversupplied. We are reducing capital by nearly 20\% and production approximately 15\%. One company is not positioned to `fix the market'.'' \\
\mbox{[\data{quoted_excerpt_labels/audit_types/q381486}]}
\end{quote}

\begin{quote}
``We proactively impaired and early retired a number of our older coastal barges and tugboats. We believe this action is necessary not only to help restore an early balance into the market, but also to improve our profitability going forward. If others will follow our lead, an accelerated balance can be restored to the struggling market. We did our part to rationalize the industry's fleet and improve our profitability going forward, but we need the industry to mirror our actions.'' \\
\mbox{[\data{quoted_excerpt_labels/audit_types/q327069}]}
\end{quote}

\subsubsection{Category 3: The industry has shown or will show supply or capacity or price discipline}
Rather than expressing a need for more supply or capacity discipline, a firm may note that it has occurred or will occur and expresses support for it. This or similar language appeared in 11\% of the true positive transcripts. These transcripts note the firm’s supply or capacity cuts in the context of an industry effort.

\begin{quote}
``Market conditions remain challenging, exacerbated by the overall oversupply in the industry. We started maintenance of our facilities and adjusted our production utilization rate to 50\% in light of weak market demand. There has been consensus among the manufacturers to promote a healthier development of the industry through exercising self-discipline.'' \\
\mbox{[\data{quoted_excerpt_labels/audit_types/q396666}]}
\end{quote}

\begin{quote}
``We have seen proactive initiatives to restore the industry's healthy development. Implementing sales discipline would be fundamental to mitigating the irrational competition amid falling prices. We plan to maintain a relatively low utilization rate in 2025 until a turning point emerges in the sector.'' \\
\mbox{[\data{quoted_excerpt_labels/audit_types/q396668}]}
\end{quote}

\begin{quote}
``If we could scare people out, and they knew how much we were going to be building [then] maybe they wouldn't build theirs. We were trying to get some attention and show that, hey, everybody else should be disciplined, and we're going to come out and throttle this. This is not a cycle where you would expect to see a tremendous amount of capacity come on.'' \\
\mbox{[\data{quoted_excerpt_labels/audit_types/q196851}]}
\end{quote}

A firm may commend competitors for acting responsibly with respect to their output or prices.

\begin{quote}
``We finally have sanity back in the seaborne iron ore market. I truly commend Rio Tinto and Vale for eliminating the reckless behavior that had infected the market for a number of years. We expect this type of responsible behavior to carry on, which will continue to support strong iron ore prices.'' \\
\mbox{[\data{quoted_excerpt_labels/audit_types/q445959}]}
\end{quote}

\begin{quote}
``The entire industry moved to action and acted responsibly last year. So we were able to, for the first time, put the prices up and that's a great thing for our business and for our industry.'' \\
\mbox{[\data{quoted_excerpt_labels/audit_types/q558080}]}
\end{quote}

Based on a recent change in market structure, a firm may forecast future discipline.

\begin{quote}
``We feel that consolidation somehow will improve market discipline going forward. Given the fact that operating in such a competitive scene has led to price erosion in data service pricing, improving market discipline is a very welcome element and consequence of this transaction.'' \\
\mbox{[\data{quoted_excerpt_labels/audit_types/q309604}]}
\end{quote}

\subsubsection{Category 4: The industry needs to raise prices}
Rather than communicating that the industry needs lower supply or capacity, a firm may announce that it needs higher prices. Such was present in 11\% of the true positive transcripts.

\begin{quote}
``Tariff repair is sorely needed for return ratios to improve. We could do it, we could lead it, we could start it, but the fact is that if competition doesn't follow, then it will hurt us. All the players need it.'' \\
\mbox{[\data{quoted_excerpt_labels/audit_types/q250287}]}
\end{quote}

\begin{quote}
``I think it's not only our desire to generate that level of profitability. So the need of increasing pricing, I think, will be more an industry need. We think we've arrived at a time where it was necessary to increase a little bit our pricing.'' \\
\mbox{[\data{quoted_excerpt_labels/audit_types/q408614}]}
\end{quote}

\begin{quote}
``We have an intention to do further price increases. The entire industry has the same need. Otherwise it wouldn't work. So there is a strong need for price increases in Consumer Tissue in Europe. Many of our competitors in Europe have lower margins than we have to start with. So they are in as much need of price increase as we are or even more so.'' \\
\mbox{[\data{quoted_excerpt_labels/audit_types/q215003}]}
\end{quote}

The called-for price increases may be in the context of responding to cost increases or ending a price war.

\begin{quote}
``We as an industry need to do a better job of tying those contracts to a CPI that actually means something to us, right? We can't be the volume leader by stealing from our competition because we know what that's going to lead to. But in a market that's growing, we should get our fair share of volumes.'' \\
\mbox{[\data{quoted_excerpt_labels/audit_types/q375388}]}
\end{quote}

\begin{quote}
``Price war destroys value, and it should be reversed soon. If you look at the structural positioning of Telecom Italia, we are by far the strongest player. And therefore I think that the others are suffering as much and more than we are suffering. So it is not something that can be sustainable. This price war is unsustainable for all the players. I consider it a missing opportunity for the whole Italian market.'' \\
\mbox{[\data{quoted_excerpt_labels/audit_types/q305524}]}
\end{quote}

\subsubsection{Category 5: The firm is being or will be less aggressive in pricing and references competitors’ conduct in a manner encouraging similar conduct}
In eight transcripts, the firm announced it has raised price or will raise price and, when referencing competitors’ pricing, encouraged them to do so or commended them for having done so.

\begin{quote}
``Our strategy of increasing the prices during quarter one has paid off. We [are] going to increase prices further in second half [as] we believe that there is still room to go in that direction. We hope that the market goes in the same direction. The market is now rational, and we believe that there is room for it to keep improving.'' \\
\mbox{[\data{quoted_excerpt_labels/audit_types/q309638}]}
\end{quote}

\begin{quote}
``We don't want to adjust prices down. We are not cutting prices further. We have been able to maintain our prices. We have seen signs that some of our competitors are doing the same.'' \\
\mbox{[\data{quoted_excerpt_labels/audit_types/q223177}]}
\end{quote}

\begin{quote}
``Our primary objective will be optimizing profits. We made a conscious decision to trade off attendance [i.e., volume] for yield. We can lean even more heavily into pricing and expect this will further lower our attendance. I can congratulate SeaWorld and Cedar Fair for how they have maintained better pricing integrity than us.'' \\
\mbox{[\data{quoted_excerpt_labels/audit_types/q526439}]}
\end{quote}

The firm may recognize its role as a market leader with regards to pricing.

\begin{quote}
``We adopted a very clear strategy of price-over-volume strategy for the bromine [and] we are acting as a market leader. When there is a need, we are reducing … our volumes in order to keep the prices at the right level. We don't want prices of bromine as the leader in the market to be too high because this will start to create incentives for other producers. We believe that the current price level of the bromine is more or less optimal.'' \\
\mbox{[\data{quoted_excerpt_labels/audit_types/q631460}]}
\end{quote}

\subsubsection{Some other examples}

Some other examples include a firm putting forth an industry objective that puts more weight on price and margins than on market share and volume. If adopted by firms, it has the implication of less aggressive competition and higher prices.

\begin{quote}
``Not only Seagate but the entire industry, decide to focus more on improving the bottom line than fighting for market share. Market share is not our focus. It's taking the right profit from this product. Industry could be constrained. Our objective is to be very reasonable with the capacity we put in place.'' \\
\mbox{[\data{quoted_excerpt_labels/audit_types/q408622}]}
\end{quote}

\begin{quote}
``Value over volume. The industry is not really equipped to really capture that value proposition. We're going to have to start, as an industry … to collaborate on turning this dynamic around [by] consolidating districts, rationalizing capital expenditures, not building duplicate infrastructure in the same camps.'' \\
\mbox{[\data{quoted_excerpt_labels/audit_types/q325064}]}
\end{quote}

Finally, one transcript facilitated hub-and-spoke collusion as a manufacturer-hub sought to dampen competition among dealer-spokes.

\begin{quote}
``We will not oversupply the market [and] we hope that our dealers will follow suit and stop competing with themselves.'' \\
\mbox{[\data{quoted_excerpt_labels/audit_types/q548932}]}
\end{quote}

\subsubsection{Summary}

The pool of true positives provides encouraging evidence for an LLM being part of a screening protocol. Out of hundreds of thousands of transcripts, the LLM managed to identify a small number that contained collusive content and would be worthy of further investigation. In addition, it did so for content spanning many different forms – reducing supply, constraining capacity, raising prices, stemming price decreases – with messages that conveyed what the industry needs to do, what it did and should continue to do, and a firm acting a leader for the industry. At the same time, it has its limitations as reflected in many false positives. We next turn to examining those false positive transcripts towards understanding how the LLM was misled.
